\sectionkicker{Source-backed technical notes}
\subsection*{Model側もTool-Nativeへ — Qwen3-Max・GLM-4.7-Flash・GLM-Image: Technical Notes}
本欄は記事本文で圧縮した一次資料上の情報を、比較・再検証しやすい形で整理したものである。時系列、確認済みの主張、留意点、一次資料URLを掲載する。
\medskip
\subsection*{Theme at a glance}
\addcontentsline{toc}{subsection}{Theme at a glance}
\begin{center}
\footnotesize
\begin{tabularx}{\linewidth}{@{}p{0.20\linewidth}p{0.10\linewidth}p{0.14\linewidth}X@{}}
\toprule
資料 & 位置づけ & 種別 & 時系列 \\
\midrule
Qwen3-Max January Model Studio updates & 主要資料 & モデル更新 & 2026-01-27 (Model Studio提供開始) \\
GLM-4.7-Flash and GLM-Image & 主要資料 & モデル更新 & 2026-01-14 (モデル公開); 2026-01-19 (モデル公開) \\
SGLang v0.5.8 & 補足資料 & Framework & 2026-01-23 (プロジェクト公開) \\
\bottomrule
\end{tabularx}
\end{center}
\normalsize

\begin{technicalnote}{Qwen3-Max January Model Studio updates}{主要資料}
\begingroup
% reader-facing Technical Notes break policy
\widowpenalty=10000
\clubpenalty=10000
\displaywidowpenalty=10000
\begin{tabularx}{\linewidth}{@{}>{\bfseries}p{0.22\linewidth}X@{}}
組織 & Alibaba Cloud / Qwen \\
種別 & モデル更新 \\
時系列 & 2026-01-27 (Model Studio提供開始) \\
\end{tabularx}
\smallskip
{\bfseries 一次資料から整理したtechnical points}
\begin{itemize}[leftmargin=1.5em,itemsep=0.35em]
\item \textbf{一次情報で確認できる事実}: Alibaba Cloud Model Studioはqwen3-max-2026-01-23を1月27日のinternational availabilityとして記録し、thinking/non-thinking modeの統合とthinking中のweb search・web extraction・code interpreter利用を説明する。
\item \textbf{一次情報で確認できる事実}: 同じlifecycle pageは1月4日、tool-oriented／codingやmultimodal modelを含む複数のQwen3・Wan 2.6系について広範なglobal availabilityを記録している。
\item \textbf{Vendor claim}: lifecycle catalog内の比較的なcapability表現はvendor報告であり、本サーベイでは独立再現していない。
\end{itemize}
\begin{minipage}{\linewidth}
% reader-facing Technical Notes generic boundary/source tail group
{\bfseries 読む際の境界}
\begin{itemize}[leftmargin=1.5em,itemsep=0.35em]
\item \textbf{分析上の留意点}: Model Studioのavailability dateはregionにより異なり得るため、記載されたdeployment regionへscopeを限定する必要がある。
\item \textbf{分析上の留意点}: lifecycle pageはliving documentationであるため、現在のdefault aliasではなくdated entryを1月chronologyの基準とする。
\end{itemize}
\begin{samepage}
{\bfseries 一次資料}
\begingroup\sloppy
\Urlmuskip=0mu plus 2mu\relax
\begin{itemize}[leftmargin=1.5em,itemsep=0.25em]
\item {\scriptsize\url{https://www.alibabacloud.com/help/en/model-studio/newly-released-models}}
\end{itemize}
\endgroup
\end{samepage}
\end{minipage}
\endgroup
\end{technicalnote}
\medskip

\begin{technicalnote}{GLM-4.7-Flash and GLM-Image}{主要資料}
\begingroup
% reader-facing Technical Notes break policy
\widowpenalty=10000
\clubpenalty=10000
\displaywidowpenalty=10000
\begin{tabularx}{\linewidth}{@{}>{\bfseries}p{0.22\linewidth}X@{}}
組織 & Z.ai \\
種別 & モデル更新 \\
時系列 & 2026-01-14 (モデル公開); 2026-01-19 (モデル公開) \\
\end{tabularx}
\smallskip
{\bfseries 一次資料から整理したtechnical points}
\begin{itemize}[leftmargin=1.5em,itemsep=0.35em]
\item \textbf{一次情報で確認できる事実}: Z.ai release notesは1月14日のGLM-Imageを記録し、autoregressive semantic understandingとdiffusion-based decodingを組み合わせたcontrollable image generation modelとして説明する。
\item \textbf{一次情報で確認できる事実}: Z.ai release notesは1月19日のGLM-4.7-Flashを、coding・reasoning・generative task向けのlightweight／free-tier GLM-4.7 variantとして記録する。
\item \textbf{Vendor claim}: quality、latency、throughputや「best-in-class」といった比較表現はvendor claimである。
\end{itemize}
\begin{minipage}{\linewidth}
% reader-facing Technical Notes generic boundary/source tail group
{\bfseries 読む際の境界}
\begin{itemize}[leftmargin=1.5em,itemsep=0.35em]
\item \textbf{分析上の留意点}: 1月のZ.ai model lineとして二つのreleaseを一括したtaskであり、単一architectureであるかのように記述しない。
\item \textbf{分析上の留意点}: performance claimはvendor報告であり、本サーベイでは独立再現していない。
\end{itemize}
\begin{samepage}
{\bfseries 一次資料}
\begingroup\sloppy
\Urlmuskip=0mu plus 2mu\relax
\begin{itemize}[leftmargin=1.5em,itemsep=0.25em]
\item {\scriptsize\url{https://docs.z.ai/release-notes/new-released}}
\end{itemize}
\endgroup
\end{samepage}
\end{minipage}
\endgroup
\end{technicalnote}
\medskip

\begin{technicalnote}{SGLang v0.5.8}{補足資料}
\begingroup
% reader-facing Technical Notes break policy
\widowpenalty=10000
\clubpenalty=10000
\displaywidowpenalty=10000
\begin{tabularx}{\linewidth}{@{}>{\bfseries}p{0.22\linewidth}X@{}}
組織 & SGLang Project \\
種別 & Framework \\
時系列 & 2026-01-23 (プロジェクト公開) \\
\end{tabularx}
\smallskip
{\bfseries 一次資料から整理したtechnical points}
\begin{itemize}[leftmargin=1.5em,itemsep=0.35em]
\item \textbf{一次情報で確認できる事実}: SGLang v0.5.8はGLM-4.7-Flashのday-zero support、diffusion serving拡張、FlashAttention 4 decoding、複数のlong-context／disaggregated-serving変更を追加した。
\item \textbf{Project claim}: 最大1.5倍のdiffusion高速化、chunked pipelineのnear-linear scaling、特定GLM4-MoE構成で65\%のTTFT改善などはprojectのrelease noteによる報告値である。
\end{itemize}
\begin{minipage}{\linewidth}
% reader-facing Technical Notes generic boundary/source tail group
{\bfseries 読む際の境界}
\begin{itemize}[leftmargin=1.5em,itemsep=0.35em]
\item \textbf{分析上の留意点}: releaseは多数のupstream PRと、異なるmodel／hardware条件のperformance claimを集約している。
\item \textbf{分析上の留意点}: model supportは実装上利用可能になったことを示すが、同等の品質やproduction validationを意味しない。
\end{itemize}
\begin{samepage}
{\bfseries 一次資料}
\begingroup\sloppy
\Urlmuskip=0mu plus 2mu\relax
\begin{itemize}[leftmargin=1.5em,itemsep=0.25em]
\item {\scriptsize\url{https://github.com/sgl-project/sglang/releases/tag/v0.5.8}}
\end{itemize}
\endgroup
\end{samepage}
\end{minipage}
\endgroup
\end{technicalnote}
\medskip
